\documentclass[11pt,a4paper]{article}
\usepackage{curso}
\begin{document}
\portada{Lección}{6.1}{Sistemas lineales desde el primer momento de un OU}

\noindent\textbf{Unidad 6.}\quad\textbf{Objetivos:} codificar $Ax=b$ como potencial cuadrático; ejecutar el protocolo de sistemas lineales (equilibrar y
promediar); relacionar el error con la temperatura y la ventana de promedio.
\textbf{Prerrequisitos:} 2.6, 5.4.

\section{Codificación}
Sea $A\in\R^{d\times d}$ simétrica definida positiva (SPD). Con
\[U(x)=\tfrac12x^\top Ax-b^\top x,\qquad\nabla U=Ax-b,\]
el único mínimo es $x^\ast=A^{-1}b$. Físicamente: osciladores armónicos acoplados; $A$ son los resortes y $b$ las fuerzas. Si $A$ no es SPD se resuelve
$A^\top Ax=A^\top b$, lo que eleva al cuadrado el número de condición $\kappa=\lambda_{\max}/\lambda_{\min}$.

\section{Dinámica y equilibrio}
Langevin sobreamortiguada (proceso de Ornstein--Uhlenbeck), con amortiguamiento $\gamma$:
\[dx=-\tfrac1\gamma(Ax-b)\,dt+\sqrt{2/(\gamma\beta)}\,dW\quad\Longrightarrow\quad x_\infty\sim\N\big(A^{-1}b,\ \beta^{-1}A^{-1}\big).\]
La solución es el \emph{primer momento} de la distribución estacionaria. El ruido garantiza ergodicidad: el promedio temporal de una trayectoria converge al
promedio del ensamble.

\begin{definicion}[ (protocolo de sistemas lineales, Aifer et al.)]
(1) Cargar $A$ y $b$. (2) Equilibrar durante $t_0\gtrsim\kappa\tau_r\ln(\kappa/\varepsilon_{\mu_0})$, con $\tau_r=\gamma/\|A\|$. (3) Promediar con un integrador
analógico $\bar x=\tau^{-1}\int_{t_0}^{t_0+\tau}x(t)\,dt$, con
\[\tau=\frac{2\kappa^2d\,\|A\|}{\beta\|b\|^2\varepsilon^2(1-P_\varepsilon)}\,\tau_r,\]
lo que garantiza $\|A\bar x-b\|/\|b\|\le\varepsilon$ con probabilidad $\ge P_\varepsilon$.
\end{definicion}

\section{Simulación \evS}
$d=50$, $A=GG^\top/(2d)$ con $G\in\R^{d\times2d}$ normal estándar (semilla 0, $\kappa\approx34$), $\gamma=1$, Euler--Maruyama con $dt=0{,}01$, $t_0=100$.
\begin{center}\begin{tabular}{lccc}\toprule
error relativo $\|\bar x-A^{-1}b\|/\|A^{-1}b\|$ & $\tau=10$ & $\tau=100$ & $\tau=900$\\\midrule
$k_BT=0{,}01$ & 6,7\,\% & 2,1\,\% & 0,84\,\%\\
$k_BT=0{,}001$ & 2,1\,\% & 0,66\,\% & 0,27\,\%\\\bottomrule\end{tabular}\end{center}
Multiplicar $\tau$ o $\beta$ por 10 divide el error por $\approx\sqrt{10}$: $\varepsilon\propto1/\sqrt{\beta\tau}$.

\section{Complejidad}
Termodinámico (sobreamortiguado): $O(d\kappa^2\varepsilon^{-2})$; subamortiguado: $O(\gamma d\kappa^{3/2}\varepsilon^{-2})$. Digital: $O(\min\{d^\omega,d^2\sqrt\kappa\})$
(gradiente conjugado o multiplicación rápida, $\omega\approx2{,}3$). La ventaja aparece para $d$ grande; el costo de carga $O(d^2)$ se estudia en 6.5.

\begin{aplicacion}
Soluciones aproximadas rápidas: precondicionadores, pasos de Newton, gradiente natural (Unidad 8), regresión con procesos gaussianos. Laboratorio L9:
comparación con gradiente conjugado y Cholesky.
\end{aplicacion}

\begin{error}
\begin{itemize}[nosep]
\item Promediar desde $t=0$ (sesgo del transitorio). \quad Usar $A$ no SPD sin transformar.
\item Esperar alta precisión: $\varepsilon\propto1/\sqrt{\beta\tau}$. \quad Ignorar el costo de carga.
\end{itemize}
\end{error}

\begin{thebibliography}{9}\small
\bibitem[30]{tla} M. Aifer et al., Thermodynamic linear algebra, \emph{npj Unconv. Comput.} 1, 13 (2024).
\bibitem[25]{mel} D. Melanson et al., Thermodynamic computing system for AI applications, \emph{Nat. Commun.} 16, 3757 (2025).
\bibitem[28]{coles} P. J. Coles et al., Thermodynamic AI and the fluctuation frontier, 2023.
\end{thebibliography}
\end{document}
